% !TeX root = ../../surface-lagrangian-torus-fibration.tex
% 2026-11-22 / v1.2: 共通 normal moment に相対的な fiber partition の変更と基底値の安定性．
\section{旧 softmax 値に適合する全域基底の parameter 化}
\label{ke:global-base-compatibility-section}

定理\ref{ke:actual-global-fibration}で得た actual fiber partition の quotient を
$Q_{j,t}$ とし，まだ $B_d(T)$ との同相写像を選ばずに
$\pi^0_{j,t}:X_t\to Q_{j,t}$ と書く．
以下の操作はこの fiber partition を変えず，quotient の parameter 化だけを選ぶ．
十分大きい固定 $j$ に対して
\[
 m=6b_j,\qquad N=N_j,\qquad \delta=4\pi/N,
 \qquad L=-\log|t|^2,\quad R=\log L
\]
とする．各 component の normal face の anchor level を同じ $q_a=\delta$ に取り，
相対葉族の補間支持を $q_a\in[\delta/2,2\delta]$ の中に取る．
Core filling を少し外まで延ばした場合にも，基底の分割にはこの anchor level を使える．
以下の $K_{i,j,t}$ と $A_{t,J}$ はこの共通 anchor で切った core と保持集合を指す．
実際の接着に平行 link を使った場合は，その間の共通 collar で旧 corner tori が
保たれるので，基底の分割をこの anchor へ取り直す．変更領域は同じ固定 component
compact 内にあり，既存の深度上界と費用評価は定数を増やすだけで保たれる．
補題\ref{ke:anchor-trimmed-neck-domain}により，この分割後の保持 neck の
moment 範囲は正確に $[-\delta,\delta]$ である．その外側の旧 neck 尾部は
対応する component core に入り，平行 link との差は変更 collar に含める．
既存の選択順序の中で $N$ を大きくし，最後に $t$ を小さくすることを許す．

三重点 $J$ の旧 corner と旧 projection を $W_{t,J}$，$P^{\rm old}_{t,J}$ とし，
その保持部分を $A_{t,J}$ とする．その quotient の旧 target 像を
\[
 D_{t,J}:=P^{\rm old}_{t,J}(A_{t,J})\subset\mathcal C_J
\]
と書く．補題\ref{ke:removed-corner-target-cap}の誤差より大きい
\begin{equation}\label{ke:eq:base-compatibility-small-scale}
 \eta=\eta_{j,t}=A\frac{1+\log(N+1)}R<\frac1{100}
\end{equation}
を取り，$A$ は固定した次数・族・chart のみに依存する十分大きい定数とする．

\subsection{全 ends に共通な normal longitude の定量的構成}

Target の二 label 部分を長さ $2$ の区間として
\[
 E_{ij}(u)=
 \begin{cases}
  e_i+u e_j,&0\le u\le1,\\
  (2-u)e_i+e_j,&1\le u\le2
 \end{cases}
\]
と表す．Normal moment の向きは，component $i$ の face が $\mu=-\delta$，
component $j$ の face が $\mu=\delta$ となるように一度だけ固定する．
この規約では generator は triple lattice の $(1,-1,0)$ に対応する．

\begin{lemma}[旧 corner 辺での共通 model profile]
\label{ke:common-edge-moment-profile}
同じ pair $ij$ に接する各三重点 $J=\{i,j,k\}$ について，旧 softmax 辺
$w_k=0$ のうち $\mu\in[-\delta,\delta]$ の部分は，$\mu$ で一対一に
parameter 化される．その旧辺 longitude を $u_{J,t}(\mu)$ と書く．
Complete simplex model のみによる一つの strictly increasing profile
$u_{0,t}(\mu)$ があり，全 ends に一様に
\begin{equation}\label{ke:eq:common-edge-profile-error}
 \sup_{|\mu|\le\delta}|u_{J,t}(\mu)-u_{0,t}(\mu)|
 \le\frac C{mR}.
\end{equation}
その endpoints は
\[
 u_{J,t}(-\delta)=\frac{\log N+O(1)}R,
 \qquad 2-u_{J,t}(\delta)=\frac{\log N+O(1)}R
\]
を満たし，model profile についても同様である．
特に $u_{0,t}(0)=1$ である．
\end{lemma}

\begin{proof}
$r_a=\ell_a-m$，$S=L-3m$ と置く．旧 softmax 辺は
$r_k=\max(r_i,r_j)/L$ である．その全体を
$v=\log(r_j/r_i)\in[-R,R]$ により
\begin{equation}\label{ke:eq:hexagon-edge-profile}
 r_i=\frac{S}{1+e^v+\max(1,e^v)/L},\quad
 r_j=e^v r_i,\quad r_k=\frac{\max(r_i,r_j)}L
\end{equation}
と明示できる．旧辺 longitude は正確に $u=1+v/R$ である．
$B=\ell_k=m+O(1)$ と書くと，$v=0$ の両側で
\[
 \frac{dB}{d\ell_i}=
 \begin{cases}1/L,&r_i\ge r_j,\\-1/(L+1),&r_j\ge r_i\end{cases}
\]
が成立する．以下の微分は $v=0$ で片側微分として読み，そこで連続につなぐ．

Normal generator を $V$，$\mathfrak a=|V|_h$ とすると，normal 深度
$a=\ell_i,c=\ell_j\gg m$ 上で
\[
 \mathfrak a\asymp a^{-1}+c^{-1},\qquad
 |\nabla_h^rV|_h\le C_r\mathfrak a
\]
である．$\partial_v$ の normal 成分の長さは上下から正定数で抑えられる．
その residual 成分の長さは $C/m$ 以下である．実際，
$|da/dv|\asymp\mathfrak a^{-1}$ と上の $dB/da$ から
\[
 \frac{|dB/dv|}B\le\frac C{\mathfrak a BL}\le\frac C m .
\]
Model moment を $\nu_0$，実際の全平均 moment を $\nu_J$ とする．
Tame 下界と
$d\nu(V')=-\kappa(V,V')$ により，向きを上のように取ると
\begin{equation}\label{ke:eq:edge-moment-weighted-derivative}
 c\mathfrak a\le\frac{d\nu_0}{dv}\le C\mathfrak a,
 \qquad
 c\mathfrak a\le\frac{d\nu_J}{dv}\le C\mathfrak a
\end{equation}
となる．Residual 成分の寄与は normal 成分の正の主項に対し $O(m^{-1})$ なので，
大きい $j$ で吸収できる．$\log\mathfrak a$ の $v$ 微分も有界である．

ここで加法定数を無視してはいけない．
補題\ref{ke:component-meridian-normalization}および
系\ref{ke:transition-normal-primitive}によって，この $\nu_J$ は全 ends で
同じ canonical normal moment $\mu$ である．
命題\ref{ke:relative-corner-patch}の証明で用いた full-triple comparison と
相対座標接続の primitive 評価，および系\ref{ke:relative-corner-patch-c2}から
\[
 \|G^*\lambda_{\log,t}-\lambda_{0,L}\|_{C^2(h)}\le C/m.
\]
Model primitive は全平均で変わらないので，generator の重みを保持して評価すると
\begin{equation}\label{ke:eq:edge-moment-weighted-error}
 |\nu_J-\nu_0|+
 \left|\frac d{dv}(\nu_J-\nu_0)\right|
 \le\frac C m\,\mathfrak a .
\end{equation}
全三重点で同じ $L,m$ と同じ complete simplex model を使うため，$\nu_0$ は
$k$ に依存しない．非正則な補間帯でも source primitive は $G^*\lambda$ である．

Dominant-product moment 評価により，$\mu=\pm\delta$ の model levels は
片方の normal 深度が $N$ と固定比の間，他方が $L-O(N+m)$ の領域にある．
$N\gg n_j\gg m$ とし，その後 $L/N$ を大きく選べば，両端とその固定相対幅の
buffer は全て上の estimates の範囲に入る．単調性により，その間の
$\mu\in[-\delta,\delta]$ の levels も同じ normal 下界を持つ．

$\nu_0(v_0)=\mu$ とする．$v_0\pm K/m$ において
\eqref{ke:eq:edge-moment-weighted-derivative}を積分し，
$\mathfrak a$ の比の有界性と\eqref{ke:eq:edge-moment-weighted-error}を用いる．
$K$ を固定して十分大きくすると，$\nu_J(v_0-K/m)<\mu<\nu_J(v_0+K/m)$ となる．
従って唯一の actual solution $v_J(\mu)$ は $|v_J-v_0|\le C/m$ を満たす．
これを $R$ で割れば\eqref{ke:eq:common-edge-profile-error}である．
端では小さい normal 深度が $N(1+O(m^{-1})+o_{t\to0}(1))$，大きい方が
$L-O(N+m)$ なので，endpoint の式も従う．
Model の $i,j$ 交換対称性により $\nu_0(0)=0$ であり，最後の主張を得る．
\end{proof}

\begin{corollary}[共通 endpoints を持つ全枝共通の profile]
\label{ke:common-neck-longitude}
$\kappa=\eta$ とする．前補題の profiles の endpoints を，target の
$4\eta$-vertex neighborhoods 内だけの increasing piecewise smooth rescaling に
よって揃えられる．その結果，
\[
 \sigma_t:[-\delta,\delta]\longrightarrow[\kappa,2-\kappa],\qquad
 \widetilde u_{J,t}:[-\delta,\delta]\longrightarrow[\kappa,2-\kappa]
\]
という increasing homeomorphisms を取り，
\[
 \sigma_t(0)=1,\qquad
 \sup_\mu|\widetilde u_{J,t}(\mu)-\sigma_t(\mu)|\le C/(mR)
\]
とできる．さらに $\mu\in[-\delta,-\delta/2]$ では両 profile は $C\eta$ 以下，
$\mu\in[\delta/2,\delta]$ では $2$ との差が $C\eta$ 以下である．
\end{corollary}

\begin{proof}
$\alpha_-=u(-\delta)$，$\alpha_+=2-u(\delta)$ と置く．
$[0,2]$ の increasing piecewise linear homeomorphism $\psi$ を，
\[
 0\mapsto0,\quad\alpha_-\mapsto\kappa,\quad
 4\eta\mapsto4\eta,\quad2-4\eta\mapsto2-4\eta,
 \quad2-\alpha_+\mapsto2-\kappa,\quad2\mapsto2
\]
で定め，中央区間では恒等とする．$\alpha_\pm\asymp\log N/R$，
$\eta=A(1+\log(N+1))/R$ だから，slopes と endpoints に関するこの式の
Lipschitz constants は一様に抑えられる．
Model と各 end にこの規則を適用し，$\sigma_t=\psi_0\circ u_{0,t}$，
$\widetilde u_{J,t}=\psi_J\circ u_{J,t}$ とする．
前補題の endpoint 誤差を含めて評価すれば表示された誤差が従う．
$\mu\in[-\delta,-\delta/2]$ の独立深度は $O(N)$ なので，元の longitude は
$O((1+\log N)/R)$ である．Rescaling の支持も同じ尺度だから最後の主張を得る．
\end{proof}

\subsection{実際の cap 境界と保持 corner の disk 構造}

\begin{lemma}[Moment cut の境界は一本の proper arc である]
\label{ke:corner-cap-disk}
上の $N$ と $t$ の選択をさらに強くすれば，各 $J=\{i,a,c\}$ について
core $i$ の quotient 境界と旧 corner の交わりは，$e_i$ に接する二つの旧境界辺を
結ぶ一本の proper embedded arc である．除かれる部分はその arc と $e_i$ 側の
二つの境界辺で囲まれた disk であり，$e_i$ を含む．
異なる三つの core の caps は互いに素である．従って $D_{t,J}$ は閉 disk であり，
その境界は三本の cap arcs と三本の旧境界 arcs が交互に並ぶ．

さらに $D_{t,J}$ から $\mathcal C_J$ への homeomorphism $r_{t,J}$ を，
旧 target からの変位 $C\eta$ 以下となるように取れる．
この写像は各 vertex の $[0,4\eta]^2$ の外で恒等であり，cap arc を
その vertex に接する二つの境界辺の長さ $\kappa$ の部分へ写す．
旧境界 arc での longitude は前系の $\psi_J$ による rescaling と一致する．
\end{lemma}

\begin{proof}
小さい target cap に含まれることだけでは，前半の位相的主張は従わない．
ここでは補題\ref{ke:closed-component-link}の境界の明示式を使う．
Core $i$ 側の source で独立深度を $a,c$ とすると，boundary の二つの normal
faces は $q_a=\delta$，$q_c=\delta$ であり，両者は monotone smooth minimum
$F(q_a-\delta,q_c-\delta)=0$ によって接続される．
Joint moment map の derivative の対称部分は負定値であるため，
$\partial_aq_a<0$，$\det Dq>0$ である．
従って $q_a=\delta$ 上では
\[
 \frac{dq_c}{dc}\bigg|_{q_a=\delta}
 =\frac{\det Dq}{\partial_aq_a}<0.
\]
Normal level は $a\asymp N$ の唯一の graph であり，$c$ が旧 boundary の
$m+O(1)$ から $2N$ へ増える間に，$q_c$ は $\delta$ より十分大きい値から
小さい値へ移る．従って二つの faces は唯一の joint-moment 近傍で接続する．
Smooth minimum の zero set はその moment coordinates で一本の滑らかな arc であり，
余分な component は生じない．

旧 boundary が $6.2b_j$ より浅い部分では，共通 plateau 上の neck face を用いる．
そこで同じ normal moment cut と厳密に一致するので，上の arc の続きを与える．
旧 softmax 境界との交点の一意性は，補題\ref{ke:common-edge-moment-profile}の
strict monotonicity による．その交点はそれぞれ $\mu=\pm\delta$ の指定された端で
ある．Fixed compact core の深度上界により，当該 corner の交点は component $i$ の
dominant chart に含まれる．この chart では他の boundary divisor の defining
functions は非零 unit であり，その $N$-level faces は入ってこない．
従って当該 chart 内での boundary は，いま列挙した二つの normal faces と一つの
rounding で尽くされる．これが余分な boundary component を排除する理由である．
旧 $e_i$ の source では二つの独立深度は $m+O(1)\ll N$ であり，これは
link の内側にある．従って除かれる側は $e_i$ を含む disk である．
三つの caps は\ref{ke:removed-corner-target-cap}の互いに素な vertex neighborhoods
に含まれるので，残る $D_{t,J}$ は表示された六つの boundary arcs を持つ disk である．

変位を制御した homeomorphism を具体的に作る．
$e_i$ を原点とする square の二座標を用い，$U_i=[0,4\eta]^2$ とする．
$A$ を大きく選んだので cap は $[0,\eta/2]^2$ の中にある．
$D_{t,J}\cap U_i$ は一つの corner を上の arc で切った閉 disk である．
その境界上で，外側の二辺は固定し，旧 coordinate-axis の二部分は
前系の $\psi_J$ によって parameter 化する．残る cap arc を，
$(\kappa,0)$ から原点を通り $(0,\kappa)$ へ至る二本の axis segments へ
orientation を保って写す．これで boundary circles の homeomorphism を得る．

ここで disk extension を仮定として残す必要はない．両方の curvilinear disks を
有限個の boundary arcs に適合する閉 disk 座標で表し，boundary circle の写像を
$e^{i\theta}\mapsto e^{i\varphi(\theta)}$ と lift する．
$\varphi(\theta+2\pi)=\varphi(\theta)+2\pi$ で狭義単調である．
\[
 (r,e^{i\theta})\longmapsto
 \left(r,e^{i[\theta+\chi(r)(\varphi(\theta)-\theta)]}\right)
\]
を用いれば，$\chi=0$ near $0$，$\chi=1$ near $1$ として boundary map を disk
全体へ延長できる．各 $r$ で angular map は狭義単調だから homeomorphism である．
ここで用いる disk 構造は上で確認した一本の proper arc による切断から得られるもの
であり，未知の global extension property ではない．有限の piecewise smooth
refinement を使えば regular strata 上の smoothness も保てる．

この extension の source と image はともに $U_i$ 内にあるので，変位は
$4\sqrt2\eta$ 以下である．三つの $U_i$ は互いに素で，外側の二辺を固定したため
その外の恒等写像と貼り合わさる．これが $r_{t,J}$ である．
\end{proof}

\subsection{全域 parameter 化の構成と旧 corner 値の近似}

\begin{proposition}[旧 softmax 値に近い全域基底値]
\label{ke:global-base-parameterization}
固定した marked tree system $T$ と map $p_T:B_d(T)\to\mathcal C_d$ を用いる．
上の $j,N,t$ の選択の下で，labels と core vertices を保つ homeomorphism
\[
 h_{j,t}:Q_{j,t}\longrightarrow B_d(T)
\]
を取り，$\pi_{j,t}=h_{j,t}\pi^0_{j,t}$，$f_{j,t}=p_T\pi_{j,t}$ とすると，次を
同時に満たすようにできる．
\begin{enumerate}
\item 全ての保持 corner について
\begin{equation}\label{ke:eq:global-base-old-corner-error}
 \sup_{x\in A_{t,J}}d_\square
 \bigl(f_{j,t}(x),P^{\rm old}_{t,J}(x)\bigr)\le C\eta_{j,t}.
\end{equation}
\item Component $i$ の core および相対葉族を変更した collar の像は
\[
 f_{j,t}(K_{i,j,t}\cup\text{変更 collar})
 \subset B_{d_\square}(e_i,C\eta_{j,t})
\]
に含まれる．
\item 三重点に向かう各 neck leaf arm の外側部分は，その end の旧 softmax
辺と共通の normal moment を用いて parameter 化される．その longitudinal
coordinate と，全枝共通の $\sigma_t(\mu)$ の差は $C\eta_{j,t}$ 以下である．
\end{enumerate}
Fiber partition，actual Lagrangian 性，連結性，source・branch・core の fiber 型は
変わらない．Parameter 化を有限の piecewise smooth strata に適合させられる．
\end{proposition}

\begin{proof}
\emph{Target triangle の明示的な inset と panels．}
各 $J$ で $\rho_J(w)=w/\sum_{a\in J}w_a$ とし，$0<\varepsilon=\eta<1/100$ に対して
\[
 H_s(w)=\rho_J^{-1}\bigl((1-3s\varepsilon)\rho_J(w)
                   +s\varepsilon(1,1,1)\bigr),\qquad0\le s\le1
\]
と置く．$H_1$ は $\mathcal C_J$ をその内部の閉 disk へ写す embedding であり，
$\sup_wd_J(H_s(w),w)\le C\varepsilon$ である．
この評価は表示された path の piecewise smooth length を積分すれば従う．

外周と inset disk の間の closed annulus は，正確に
$H_s(w)$，$w\in\partial\mathcal C_J$，$0\le s\le1$ で parameter 化される．
実際，simplex 座標 $\lambda$ では
$s=\min_a\lambda_a/\varepsilon$ であり，元の boundary point は
$(\lambda-s\varepsilon\mathbf1)/(1-3s\varepsilon)$ で回復できる．
従って annulus の存在やその trivialization を追加仮定していない．

Outer boundary を六つの arcs に分ける．Vertex $e_i$ の arc $I_i$ は，隣接する
二つの辺上で $e_i$ からの長さが $\kappa$ 以下の部分の和とする．
残る三つの edge arcs は $I_{ij}=E_{ij}([\kappa,2-\kappa])$ である．
これらの arcs と $s\in[0,1]$ の積の像を target panels とする．
Vertex panel $H([0,1]\times I_i)$ は
$B_{d_J}(e_i,C(\kappa+\varepsilon))$ に含まれる．

\emph{保持 corner 全体を中央 disk として使う．}
$Q_{j,t}$ 内の $A_{t,J}$ の quotient は $D_{t,J}$ と同一視できる．
この\emph{全体}に対して
\[
 g_{J,t}=H_1\circ r_{t,J}
\]
を用いる．補題\ref{ke:corner-cap-disk}により，これは inset disk への
homeomorphism であり，旧 $P^{\rm old}$ との差は $C\eta$ 以下である．
保持 corner をさらに狭い $B=6.5b_j$ の帯などへ切り直してはいないので，
\eqref{ke:eq:global-base-old-corner-error}は全保持点で成立する．

\emph{Actual quotient の残りの pieces．}
一つの triangle cell $\tau_J$ から $D_{t,J}$ の内部を除いた領域は，三つの core
sectors と三つの neck leaf-arm rectangles に分かれる．
Core sector は core cone の link graph における該当 edge の cone である．
その境界 link edge は，二つの branch vertices からの neck tails と，中央の cap arc
からなる．Leaf-arm rectangle は，pair $ij$ の marked tree の葉 $k$ に向かう arm
を，branch vertex と旧 softmax 辺の間で切ったものに対応する．

この分割を actual quotient 内で行える理由を確認する．
旧辺 $w_k=0$ は $B=m+O(1)$ の graph で，normal 深度は
$\mu\in[-\delta,\delta]$ に一様に $cN$ 以上である．
$N\gg n_j$ なので，この graph の全 fibers は K2 の共通 plateau overlap 内にある．
そこで normal moment は同じ $\mu$ で，reduced 葉は順序を保った同じ全 tori である．
この marked Reeb continuation は，既存の diffeomorphism 族から得られる．
補題\ref{ke:fixed-marked-tree-exhaustion}の $D_{j,\mu,t}\circ S_t$ と，
定理\ref{ke:saddle-moving-relative-interpolation}の $h_\mu^{\chi(\mu)}$
（各補間 band の外では恒等）を合成し，固定した基準葉族
$\mathcal F^0$ を構成後の $\mathcal F_\mu$ へ送る族 $\Xi_{\mu,t}$ を取る．
系\ref{ke:actual-saddle-moving-neck-interpolation}はこの葉族の全逆像を使っている．
従って商上の写像
\[
 (\mu,[p]_{\mathcal F_\mu})
 \longmapsto(\mu,[\Xi_{\mu,t}^{-1}(p)]_{\mathcal F^0})
\]
は $\mu$ と labels を保つ graph の自明化である．
各 leaf arm 上で branch vertex から boundary graph までの正の長さを
再規格化すれば，arm は $[-\delta,\delta]\times[0,1]$ と parameter 化できる．
最後の座標 $r$ は branch で $0$，旧 corner boundary で $1$ となる．
各 $\mu$ で端の長さを再 parameter 化するだけであり，foliation の変更ではない．
Source がこの rectangle 内にあっても quotient は同じ product であり，対応する
一つの base value の fiber が pinched torus になるだけである．

\emph{Leaf-arm rectangles の具体的な写像．}
該当する rectangle に対して
\begin{equation}\label{ke:eq:leaf-arm-base-interpolation}
 U(\mu,r)=(1-r)\sigma_t(\mu)+r\widetilde u_{J,t}(\mu),\qquad
 F_J(\mu,r)=H_r(E_{ij}(U(\mu,r)))
\end{equation}
と定める．各 $r$ で $U$ は $[-\delta,\delta]$ から
$[\kappa,2-\kappa]$ への increasing homeomorphism である．
従ってこれは rectangle から該当 edge panel への homeomorphism である．
$r=1$ では $g_{J,t}$ と一致し，$r=0$ では $E_{ij}(\sigma_t(\mu))$ となる．
$\mu=\pm\delta$ では $U$ がそれぞれ $\kappa,2-\kappa$ に等しいため，
core sector との二本の共通 sides も既に指定される．
全 ends の longitude の共通選択と endpoint の一致は前系が保証する．

\emph{Core sectors の extension．}
Component $i$ の core cone の radial coordinate を $r_i\in[0,1]$ とする．
各 branch ray 上では，$i$ から出発する target edge の longitude を
$\kappa r_i$ と指定する．Core sector の残りの境界では，二つの neck tails に
\eqref{ke:eq:leaf-arm-base-interpolation}の $\mu=\pm\delta$ での値を，cap arc に
$g_{J,t}$ の値を用いる．これらは core sector の boundary circle を vertex panel
の boundary circle へ orientation を保って一対一に写す．
補題\ref{ke:corner-cap-disk}の証明に記した angular lift の式を用いて，この境界写像
を disk 全体へ延長する．従って core sector 全体の像は該当 vertex panel に入り，
その全点は $e_i$ から $C\eta$ 以下である．

\emph{内部木と bigons．}
Pair $ij$ に属する全 branch edges と bigons の和集合は
\[
 \bigl(T_{ij}^{\rm int}\times[0,1]\bigr)/
 \bigl(T_{ij}^{\rm int}\times\{0\}\sim v_i,
       T_{ij}^{\rm int}\times\{1\}\sim v_j\bigr)
\]
である．Actual quotient にも，中央の $\mu\in[-\delta,\delta]$ の product と，
両端の core cones による同じ表示がある．
中央の全内部木で longitude を共通の $\sigma_t(\mu)$ とし，両 core cones では
それぞれ $\kappa r_i$，$2-\kappa r_j$ とする．内部木の座標は変えない．
$p_T$ の定義で使う simplex-edge coordinate は
\[
 \Theta(u)=\bigl(\rho_d(E_{ij}(u))\bigr)_j
 =\begin{cases}u/(1+u),&u\le1,\\1/(3-u),&u\ge1\end{cases}
\]
だから，target の上の表示の第二因子には $\Theta(u)$ を用いる．
これにより全 bigons と branch edges の homeomorphism を得る．
その branch 上の値は，隣接する全 triangle panels で指定した値と同じである．

\emph{全域への貼合せと $p_T$ との関係．}
各 triangle では，上で構成した写像は全 $\mathcal C_J$ への homeomorphism である．
その後に $(p_T|_{\tau_J})^{-1}$ を合成し，bigons と branch edges では直前の
product-coordinate map を用いる．全共通 boundaries で一致し，core vertices を
正しく同定しているので，まず有限個の飽和閉片上で
$X_t\to B_d(T)$ を連続に貼り合わせることができる．
これは全 fiber 上で一定なので，quotient の普遍性により
連続写像 $h_{j,t}:Q_{j,t}\to B_d(T)$ を誘導する．
各 labelled cell の内部で一対一かつ全射であり，boundary の identifications も同じ
だから，この写像は一対一の連続全射である．Compact から Hausdorff への写像なので
homeomorphism である．これで $f_{j,t}=p_Th_{j,t}\pi^0_{j,t}$ が，構成中の値と
厳密に同じであることが確認できた．
系\ref{topo:cubical-composite-connected}により，この同じ $f_{j,t}$ の全 fiber は連結である．
この結論は後の距離評価又は GH 収束を入力にしていない．

Regular strata の座標には有限の piecewise smooth refinement を用いる．
これは quotient の座標だけの変更であり，全 actual fibers とその isotropic strata
を保存する．必要なら angular extension の cutoff を中心近傍で零として，
各 open two-dimensional piece 上で smooth な座標を選べる．

最後に quantitative assertions を確認する．保持 corner の誤差は
$r_{t,J}$ と $H_1$ の変位を加えれば $C\eta$ 以下である．
Core の triangle sectors は vertex panels に，bigon sectors と branch rays は
$e_i$ から長さ $\kappa$ 以下の edge parts に入る．
変更 collar の外側部分では，normal moment が $\delta$ と固定比の間にあるため，
前系の最後の評価から $\sigma_t$ と $\widetilde u_{J,t}$ は該当する endpoint の
$C\eta$ 近傍にある．\eqref{ke:eq:leaf-arm-base-interpolation}もその近傍を保ち，
$H_r$ による追加変位は $C\eta$ 以下である．Corner の変更 collar は既存の target
cap estimate と保持 corner の比較で扱える．これで全ての主張を得る．
\end{proof}

\subsection{内部木の共通 moment と neck 全体の値}

\begin{proposition}[潰される内部木上では共通 normal moment を保てる]
\label{ke:collapsed-tree-common-moment}
命題\ref{ke:global-base-parameterization}の parameter 化を用いる．
$u\in(\kappa,2-\kappa)$ とし，$\mu=\sigma_t^{-1}(u)$ と置く．
このとき
\[
 f_{j,t}^{-1}(E_{ij}(u))
\]
は，pair $ij$ の neck において，同じ normal moment $\mu$ 上で
$T_{ij}^{\rm int}$ 全体に属する構成後の neck fibers を合併した集合に等しい．
相対補間の支持外では，これらは旧 K2 fibers そのものである．
Core に近い $u\in[0,\kappa]$ 又は $u\in[2-\kappa,2]$ の部分は，それぞれ
component $i$ 又は $j$ の固定 compact core と変更 collar に含まれる．
前者の厳密な同一 moment の主張は，corner 値の $O(\eta)$ 近似とは別の性質である．
\end{proposition}

\begin{proof}
$p_T$ の $E_{ij}(u)$ 上の逆像は，定義通り一つの $T_{ij}^{\rm int}$ である．
Triangle の内部は三 label の support を持つ target 内部へ homeomorphically 写るので，
そこから追加の fiber は来ない．中央の book の全内部木では longitude を
$\sigma_t(\mu)$ と共通に選んだため，その逆像は同じ $\mu$ の slice に厳密に一致する．
両端の longitude ranges には各 component の core cones を対応させたので，残りの
主張も従う．変更 collar の境界を別の平行 level に取った場合は，例外をその
component の固定 compact collar に拡大すれば同じ結論を使える．
\end{proof}

\begin{corollary}[旧 neck と source critical values の target control]
\label{ke:neck-base-value-control}
保持された pair $ij$ の neck 上で，normal moment を $\mu$ とすると，
\[
 d_\square\bigl(f_{j,t}(x),E_{ij}(\sigma_t(\mu))\bigr)\le C\eta_{j,t}.
\]
Core collars の外にある内部木の部分では左辺は零である．
特に全 source critical values の像は
$B_{d_\square}(e_i+e_j,C\eta_{j,t})$ に含まれる．
\end{corollary}

\begin{proof}
内部木では前命題を使う．旧 corner の外の leaf-arm rectangle では
\eqref{ke:eq:leaf-arm-base-interpolation}と
$|\widetilde u_{J,t}-\sigma_t|\le C/(mR)$ を使えばよい．
残る old neck--corner overlap では $B\in[m+O(1),7b_j]$ である．
固定 $\mu$ の source level を $a=a(B)$ と書くと，weighted normal estimates より
\[
 |\partial_B\nu|\le C\mathfrak a/B,
 \qquad -\partial_a\nu\ge c\mathfrak a^2.
\]
従って $|da/d\log B|\le C/\mathfrak a$ であり，
\[
 \left|\frac{d}{d\log B}\log\frac{r_j}{r_i}\right|\le C.
\]
この $B$-interval の ratio は有界なので，旧 boundary $w_k=0$ と同じ $\mu$ で比較
した longitude の差は $C/R$ 以下である．また当該 overlap の第三座標は
\[
 0\le w_k=1+R^{-1}\log(r_k/r_{\max})
 \le C(1+\log b_j)/R\le C\eta_{j,t}.
\]
第三座標を零にする線分と，辺に沿う longitude の差を合わせ，保持 corner 上の
\eqref{ke:eq:global-base-old-corner-error}を用いれば求める評価を得る．
Source の critical moment は $\mu=0$ であり，$\sigma_t(0)=1$，
$E_{ij}(1)=e_i+e_j$ だから最後の主張が従う．
\end{proof}

\paragraph{厳密一致の障害と，得られた結論の範囲．}
旧 softmax 辺 $w_k=0$ の torus は actual quotient では triangle の regular
interior にあり，対応する cellular boundary は branch locus である．
$p_T$ は triangle interior を target triangle interior へ写すので，保持 corner
全体で $p_T\pi=P^{\rm old}$ とすることは一般に不可能である．
上の inset はこの障害を $O(\eta)$ の誤差で解消する．

既存の slow diagonal によって $\log N_{j(t)}=o(R_t)$ とすれば
$\eta_{j(t),t}\to0$ である．従って保持 corner の旧 target 値，component の小さい
target caps，および潰される内部木の共通 moment は，全域の一つの
$f_t=p_T\pi_t$ に対して同時に実現される．
本節は target の parameter 化を構成したものであり，normal levels 又はその全木の
actual ambient 径，全曲面経路に対する no-shortcut 下界を証明したものではない．

\subsection{Normal moment を保つ相対変更の安定性}
\label{ke:relative-replacement-section}

作用 atlas のために reduced 葉や core filling を選び直す場合，
それらを結ぶ diffeomorphism の微分を actual KE 計量で小さく抑える必要はない．
以下は，変更後の局所構成とその全 fiber 接着が別途得られた場合に，
同じ coarse target 上の値を制御する命題である．
この命題自体は，critical actions を指定する局所構成の存在を与えない．

\begin{proposition}[共通 moment に相対的な変更の基底値比較]
\label{ke:moment-relative-replacement}
命題\ref{ke:global-base-parameterization}の許容される固定 $j,t$ と，
同じ anchor cuts，保持 corner $A_{t,J}$，component core と変更 collar，
保持 neck の有限閉被覆を用いる．この被覆上に別の局所 actual Lagrangian
fibration データを与え，次を仮定する．
\begin{enumerate}
\item 各 $A_{t,J}$ と全 K2 交わりを含む開 end collars では，
元の全 fiber partition を保つ．Source の compact nodal neighborhoods も保つ．
\item Pair $ij$ の保持 neck では，同じ actual normal circle，
向きと加法定数を固定した同じ canonical moment $\mu$，
および同じ anchor levels $\mu=\pm\delta$ を用いる．
Reduced 葉の商には，$d-2$ 個の labelled ends を持つ許容される三価の
固定 pants tree $\widehat T_{ij}$ による
実際の marked product 自明化があり，$\mu$ と全 end labels を保つ．
各 leaf arm は branch から旧 corner boundary まで
$[-\delta,\delta]\times[0,1]$ と parameter 化され，
source values を含めて連続で，正則 strata 上で piecewise smooth とする．
木の抽象的な同型だけをこの自明化の代わりには用いない．
\item 各 component の変更は同じ core と変更 collar の内部に限り，
局所商は $\widehat T=(\widehat T_{ij})$ の labelled link の cone である．
その boundary の neck tails，branch rays，および保持 corner の cap arcs は
上の marked product と一致する．特に cap arcs と anchor cuts は変更しない．
\item 各局所商は compact Hausdorff で，全交わりは双方の全 fibers の和であり，
その partition は一致する．局所 fibers は連結で，正則部分の Lagrangian 性と
特異部分の等方的 strata は各局所構成で既に確認されているものとする．
\end{enumerate}
すると，これらのデータの閉商 $\widehat Q_{j,t}$ と，
labels を保つ homeomorphism
$\widehat h_{j,t}:\widehat Q_{j,t}\to B_d(\widehat T)$ が得られる．
対応する全域 fibration を $\widehat\pi_{j,t}$，
$\widehat f_{j,t}=p_{\widehat T}\widehat\pi_{j,t}$ とすると，
parameter 化を
\begin{equation}\label{ke:eq:relative-replacement-value}
 \widehat f_{j,t}=f_{j,t}\quad\hbox{on every }A_{t,J},\qquad
 \sup_{x\in X_t}d_\square
       \bigl(\widehat f_{j,t}(x),f_{j,t}(x)\bigr)
 \le C\eta_{j,t}
\end{equation}
となるように選べる．$C$ は固定した次数・族と target の有限幾何に依存してよいが，
reduced 葉を変更する自明化や diffeomorphism の微分には依存しない．
\end{proposition}

\begin{proof}
まず補題\ref{glue:closed-whole-fiber-quotient}を適用する．
交わる local fibers は同じ全 ambient subset なので，local fiber relations の
有限和は既に推移的で，閉同値関係になる．従って compact Hausdorff な
$\widehat Q_{j,t}$ を，後の $B_d(\widehat T)$ との同定に依存せず得る．
各 global fiber は local fiber そのものであり，連結性と actual Lagrangian 性，
特異 strata の条件も保持される．

保持 corner の商には，命題\ref{ke:global-base-parameterization}で用いた
\emph{同じ} $g_{J,t}=H_1r_{t,J}$ を用いる．Cap arcs を変えていないので，
domain $D_{t,J}$ とこの写像はそのまま使える．
新しい leaf-arm rectangle の座標を $(\mu,\widehat r)$ とすると，
\[
 \widehat U(\mu,\widehat r)
 =(1-\widehat r)\sigma_t(\mu)+\widehat r\widetilde u_{J,t}(\mu),
 \qquad
 \widehat F_J(\mu,\widehat r)
 =H_{\widehat r}\bigl(E_{ij}(\widehat U(\mu,\widehat r))\bigr)
\]
と置く．各 $\widehat r$ で $\widehat U$ は同じ interval の increasing
homeomorphism である．従ってこれは同じ target panel への homeomorphism で，
$\widehat r=1$ では保持 corner，$\widehat r=0$ では
$E_{ij}(\sigma_t(\mu))$ と一致する．
ここで新しい $\widehat r$ と古い $r$ が近いことは要求しない．

内部木全体では，同じ $\sigma_t(\mu)$ と simplex-edge coordinate $\Theta$
を使う．Core cones の branch rays は $\kappa r_i$ または $2-\kappa r_i$
へ送り，triangle sectors は既に指定した boundary map を vertex panels へ
延長する．Cap arcs と全 end labels の整合性により，
命題\ref{ke:global-base-parameterization}の disk extension をそのまま使える．
新しい内部木の組合せ型が元のものと異なっても，対応する bigons と branch
edges に同じ product-coordinate construction を用いればよい．

これらの写像は有限閉 pieces 上で一致し，その全 fibers 上で一定なので，
$X_t$ 上の連続写像を先に貼り合わせ，閉商の普遍性で
$\widehat h_{j,t}$ を得る．各 labelled cell の内部で全単射であり，
boundary の identifications も一致するため，
compact-to-Hausdorff の判定により同相である．
新たな atlas が与えられている場合には，この同相でその atlas を移すだけである．
Parameter 化から critical-action の直線性を作り出す操作ではない．

値の比較は全被覆の三種類の領域ごとに行う．保持 corner では二つの写像は同じである．
一つの component の core と変更 collar では，両像とも同じ $e_i$ の
$C\eta_{j,t}$ 近傍にある．残る pair $ij$ の neck では，元の写像に
系\ref{ke:neck-base-value-control}を使う．新しい写像にも，表示した
$\widehat F_J$ と $|\widetilde u_{J,t}-\sigma_t|\le C/(mR)$ から
\[
 d_\square\bigl(\widehat f_{j,t}(x),
                    E_{ij}(\sigma_t(\mu(x)))\bigr)\le C\eta_{j,t}
\]
が成り立つ．内部木ではこの距離は零である．共通の同じ moment を使うため，
三角不等式で\eqref{ke:eq:relative-replacement-value}を得る．
この最後の評価には新旧の graph 座標の差も，その微分も入らない．
\end{proof}

\begin{corollary}[相対変更後の同じ射影による GH 実現]
\label{ke:relative-replacement-gh}
定理\ref{ke:global-metric-realization}の遅い対角選択に沿って，
前命題の局所データが一つの固定 $\widehat T$ に対して得られたとする．
すると同じ actual KE 距離について
\[
 \operatorname{dis}(\widehat f_t)
 \le\operatorname{dis}(f_t)+C\eta_{j(t),t}\longrightarrow0.
\]
各 $\widehat f_t$-fiber，すなわち一つの $p_{\widehat T}$-fiber 上の
全 $\widehat\pi_t$-fibers の合併について，
\[
 \sup_{z\in\mathcal C_d}
 \operatorname{diam}_{R_t^{-1}d_{g_t}}\widehat f_t^{-1}(z)
 \le\operatorname{dis}(\widehat f_t)\longrightarrow0
\]
である．従って個々の $\widehat\pi_t$-fibers の ambient 直径も一様に零へ向かう．
特に $d=5,6$ で前命題の局所データが作用適合 atlas を持つよう構成できれば，
その atlas と GH 実現は同じ $\widehat\pi_t$ によって達成される．
\end{corollary}

\begin{proof}
各二点の target 距離の差は\eqref{ke:eq:relative-replacement-value}の二倍以下である．
これと $\eta_{j(t),t}\to0$ を使い，定数 $C$ を増やす．
$\widehat f_t$ は全射であり，同じ target 値を持つ二点の正規化距離は
$\operatorname{dis}(\widehat f_t)$ 以下だから，fiber の主張も従う．
Atlas の存在は前命題の入力であって，この距離比較からの結論ではない．
\end{proof}

\begin{remark}[Normal moment の絶対誤差だけでは足りない]
\label{ke:raw-moment-error-insufficient}
共通の canonical moment を保つ条件を，単に
$|\widehat\mu-\mu|=o(1)$ へ弱めることはできない．
実際，固定 stage と $0<\alpha<\beta<1$ に対し，
$\mu_t=-L_t^{-\alpha}$，$\nu_t=-L_t^{-\beta}$ と取る．
十分小さい $t$ で両者は保持 moment interval に入り，
補題\ref{ke:common-profile-log-depth}より
\[
 |\mu_t-\nu_t|\longrightarrow0,\qquad
 \sigma_t(\mu_t)\longrightarrow\alpha,\qquad
 \sigma_t(\nu_t)\longrightarrow\beta.
\]
従って対応する辺上の距離は $\beta-\alpha>0$ へ収束する．
近似的な matching を使うなら，必要なのは
$|\widehat\sigma_t(\widehat\mu)-\sigma_t(\mu)|$ の一様な制御である．
この例は profile の誤差尺度を区別するものであり，
actual atlas fibration の不存在を示す反例ではない．
\end{remark}
